\section{Направленные ошибки и наблюдения с конечной дисперсией}
\label{diroracle:section}
Дадим две реализации \eqref{nloracle:query} и один способ получения
наблюдений. Медиана средних является стандартной конструкцией~\cite{dirlugosi};
робастные стохастические варианты FW изучаются
в~\cite{dirtang,dirsfyraki}. Здесь она обеспечивает сертификат
глобальной ошибки при точных значениях функции.

\begin{proposition}[Векторные наблюдения и оценки атомов]
\label{diroracle:width-transfer}\label{diroracle:scores}
Определим полунорму допустимых направлений
\[
 \omega_Q(v)=\max_{u,w\in Q}\langle v,u-w\rangle\le D\|v\|.
\]
Если векторный запрос удовлетворяет
$\omega_Q(\widehat g-\nabla f(x))\le\eta$, выберем
$s\in\arg\min_{v\in Q}\langle\widehat g,v\rangle$ и
$h=\langle\widehat g,x-s\rangle$.
Тогда выполнено \eqref{nloracle:query}; гарантия верна и для
выбранного направления away с его приближённым зазором.

Пусть вместо этого $Q=\operatorname{conv}\{v_0,\ldots,v_{n_A-1}\}$,
$n_A\ge2$, а $x=\sum_i\lambda_i v_i$ задана точными неотрицательными
весами с суммой один. Положим
$t_i=\langle\nabla f(x),v_i-v_0\rangle$.
Предположим, что $q_0=t_0=0$ и
\begin{equation}\label{diroracle:score-error}
 \max_i(q_i-t_i)-\min_i(q_i-t_i)\le\eta.
\end{equation}
Возьмём $q_x=\sum_i\lambda_iq_i$, выберем
$i_s\in\arg\min_i q_i$ и положим $h=q_x-q_{i_s}$.
Это также даёт \eqref{nloracle:query} и гарантию для направления away.
Оценки $q_i$ не обязаны быть значениями одного линейного функционала.
Каждый запрос требует полного просмотра $n_A$ оценок.
\end{proposition}
\begin{proof}
Полунорма ограничивает абсолютную ошибку на разности любых
допустимых точек величиной $\eta$.
Максимизация даёт границу зазора; подстановка выбранного направления
даёт границу наклона.
Для оценок атомов положим $e_i=q_i-t_i$,
$e_x=\sum_i\lambda_i e_i$.
Число $e_x$ лежит между наименьшим и наибольшим $e_i$, поэтому
\[
 \langle\nabla f(x),x-v_i\rangle=q_x-q_i-e_x+e_i\le h+\eta.
\]
Выбранный наклон FW не меньше $h-\eta$, а наклон away
не меньше $q_i-q_x-\eta$.
Кроме того, $h\ge0$, поскольку $q_x$ есть выпуклая комбинация оценок.
\end{proof}

При евклидовой границе ошибки $\eta_{\rm norm}$ можно взять
$\eta=D\eta_{\rm norm}$ и $P^+=\|\widehat g\|+\eta_{\rm norm}$.
Последняя граница позволяет использовать совместный сертификат.
Из одних направленных ошибок или ошибок оценок атомов граница
нормы не следует.
Для шара $\ell_1$ радиуса $R$ имеем $\omega_Q(v)=2R\|v\|_\infty$;
для симплекса она равна $\max_i v_i-\min_i v_i$.
Перебор оценок не считается одним вызовом LMO на неявно заданном многограннике.

Пусть одно новое наблюдение в запрошенной точке даёт $n_A-1$
скалярных значений $Y_i$.
Условно на всём прошлом $\mathcal F$ наблюдения внутри запроса
независимы и одинаково распределены, с известной общей границей
\begin{equation}\label{diroracle:moment}
 \mathbb E[Y_i\mid\mathcal F]=t_i(x),\qquad
 \mathbb E[(Y_i-t_i(x))^2\mid\mathcal F]\le\sigma_A^2<\infty.
\end{equation}
Координаты одного наблюдения могут зависеть друг от друга.
Пронумеруем все запросы, включая уточнения, числами $j\ge1$.
На уровне точности $r$ положим $\eta_r=\eta_0 2^{-r}$ и
\begin{equation}\label{diroracle:batch}
 b_r=\max\{1,\lceil16\sigma_A^2/\eta_r^2\rceil\},\qquad
 k_j=2\left\lceil4\log\frac{(n_A-1)j(j+1)}\zeta\right\rceil+1,
 \quad 0<\zeta<1.
\end{equation}
Возьмём $b_rk_j$ наблюдений.
Для каждой координаты вычислим медиану средних по $k_j$
непересекающимся блокам размера $b_r$ и положим $q_0=0$.

\begin{theorem}[Единое событие доверия для всех запросов]
\label{diroracle:sampling}
С вероятностью не меньше $1-\zeta$ каждый запрос алгоритма
удовлетворяет \eqref{diroracle:score-error}.
На этом событии верны все оценки теоремы~\ref{nloracle:main},
включая $\delta(x_N)\le U_N\le\varepsilon$ в случайный момент остановки.
При фактическом $T=N+1+R$ число исходных наблюдений удовлетворяет
\begin{align}\label{diroracle:samples}
 S&\le k_T\left[(N+1)b_R+\sum_{r=0}^{R-1}b_r\right]\\
  &\le k_T\left[T+16\sigma_A^2\eta_R^{-2}
       \left(N+1+\frac{1-4^{-R}}3\right)\right].\nonumber
\end{align}
Каждое наблюдение возвращает все $n_A-1$ координат, то есть
получено $(n_A-1)S$ скалярных направленных значений.
Медианы, просмотры и точные вычисления функции учитываются отдельно.
\end{theorem}
\begin{proof}
Условное неравенство Чебышёва ограничивает вероятность ошибки
одного среднего больше $\eta_r/2$ величиной $1/4$.
При $\sigma_A=0$ эта вероятность нулевая.
Независимость блоков и неравенство Хёффдинга дают вероятность
ошибки нечётной медианы не больше $e^{-k_j/8}$.
Суммирование по координатам даёт условную вероятность неудачи
не больше $\zeta/[j(j+1)]$.
Если все координатные ошибки не превосходят $\eta_r/2$,
их размах не превосходит $\eta_r$, что даёт
\eqref{diroracle:score-error}.
Применим эту оценку на событии, что запрос $j$ сделан,
возьмём ожидание и просуммируем по всем $j\ge1$.
Так как сумма $1/[j(j+1)]$ равна единице, все запросы верны
одновременно с вероятностью не меньше $1-\zeta$.
Детерминированное доказательство для каждого префикса обеспечивает
конечную остановку без заранее заданного числа запросов.

На каждом уровне $0,\ldots,R-1$ имеется один неудачный запрос,
а $N+1$ завершающих запросов имеют уровни не больше $R$.
Монотонность $k_j$ даёт первую оценку числа наблюдений.
Для второй используем $b_r\le1+16\sigma_A^2\eta_r^{-2}$ и
$\sum_{r<R}\eta_r^{-2}=\eta_R^{-2}(1-4^{-R})/3$.
\end{proof}

Конечная дисперсия направленных наблюдений возможна без конечного
момента векторного наблюдения.
На симплексе общая компонента Коши $W\mathbf1$ сокращается
в каждой разности $\langle Z,e_i-e_0\rangle$.
Сами разности всё равно должны удовлетворять условной границе
дисперсии \eqref{diroracle:moment}.
Конструкция требует конечного явного набора атомов, новых выборок,
известной границы дисперсии и точных значений функции.
